% Nouveau EGA V (cote 25): archive cote 25, site page 130 (PDF page 131).
% typescript. Draft compared: none.
% AI-transcribed by Claude Opus 5.5 (claude-opus-5-5, adaptive thinking, effort high), Message Batches API, batch ega25b (msgbatch_01A7n5YuXubUPNK1zsepi8RG), 2026-11-13; not yet checked by eye.
\scanpage{25/130}
\mspage{\hand{33}}

ce qui nous ramène au cas où S est artinien local, S=Spec(A). On procède
par récurrence sur l'entier n tel que $m^{n+1}=0$ , ce qui nous permet de
supposer que les deux sections envisagées \inserted{u,v} sont égales mod $m^{n\xout{\unsure{+1}}}$. Donc
l'une se déduit de l'autre à l'aide d'un élément \xout{xxxxxx(xxxx xxxxxxx} \hand{$\delta$ de}
$\mathrm{Hom}_k(u_o^{\hand{*}}(\hand{\Omega}^1_{F_o/k}), V)$ , où $k=A/m$ est le corps résiduel, $F_o=F\otimes_A k$ est la
fibre réduite, $V=m^n = m^n/m^{n+1}$ . Il faut prouver que \hand{$\delta$}$= 0$, moyennant
l'hypothèse $h(u)=h(v)$ . Le principe général de la vérification est le
suivant: exprimant d'abord que $h(u)$ et $h(v)$ coincident mod $m^n$, on voit
que leur ``différence'' peut s'écrire comme un élément \hand{$\delta$}$'$ de
$\mathrm{Hom}_k(w_o^{\hand{*}}(\hand{\Omega}^1_{G_o/k}),V)$ , où $w_o = h_o(u_o)=h_o(v_o)$ \xout{xxx} et où $G_o = G\otimes_A k$ ;
cet élément n'est autre que celu\uline{i} qui est déduit de \inserted[typed]{\hand{$\delta$} par composition avec} l'homomorphisme natu-%
tel $h_o^{\hand{*}} : w_o^{\hand{*}}(\hand{\Omega}^1_{G_o/k}) \hand{$\to$} u_o^{\hand{*}}(\hand{\Omega}^1_{F_o/k})$ déduit de $h_o: F_o\hand{$\to$} G_o$. Comme
$h(u)=h(v)$ donc \hand{$\delta$}$' = 0$, le composé de \hand{$\delta$} avec l'homomorphisme précédent
$h_o^{\hand{*}}$ est nul, donc si on sait que $h_o$ est surjectif, il s'ensuit que \hand{$\delta$}$=0$
et on gagne. Or le fait que $h_o:F_o\hand{$\to$} G_o$ est un monomorphisme, donc induit
un monomorphisme pour l'ensemble des points à valeurs dans les nombres
duaux sur k, implique bien en effet que $h_o^{\hand{*}}$ est surjectif (son transposé
étant injectif). Ce raisonnement est bien valable tel quel lorsque G
est représentable, ce qui n'est pas le cas cependant dansle cas qui nous
intéresse. \deleted{Il faut donc} \inserted{On peut cependant encore} définir un vectoriel \hand{$\mathcal{V}_{w_o}$} \inserted{(espace tangent à $G_o/k$ en $w_o$)} sur k, jouant le rôle
du \inserted{dual de} $w_o^{\hand{*}}(\hand{\Omega}^1_{G_o/k})$ pour exprimer la ``déviation'' de deux points de G qui
coincident mod $m^n$ \inserted{comme un élément de $\mathcal{V}_{w_o}\otimes_k V$}. C'est essentiellement straightforward, et contenu
dans les développements systématiques du par. 2\unsure{\over{5}{6}} (? Extensions infinitési-%
males), auquel on pourrait renvoyer ici. Dans le cas a), G étant le
foncteur $\mathrm{Div}_{X/S}$ , $w_o$ correspond donc à un diviseur de Cartier \xout{xx} $D_o$ sur
$X_o = X\otimes_A k$ , et il faut prendre \hand{$\Omega$}$= H^o(D_o, \underline{n}_{D_o/X_o})$ où le $\underline{n}$ est le
faisceau\xout{x} normal de $D_o$ dans $X_o$ , isomorphe également au faisceau induit
sur $D_o$

\mspage{\hand{130}}
